% tex8web: math. tex4ht's mathjax option leaves each formula's LaTeX source
% in the page, for MathJax to typeset. This file configures MathJax and
% adapts that source where MathJax differs from TeX.
\makeatletter

% the MathJax configuration, as JSON, put in the page by tex8web-page.tex:
% - extensions for the packages the document loads that MathJax has (others
%   load when a macro needs them), and always mathtools, for
%   \DeclarePairedDelimiter in macros.lua's definitions
% - no automatic equation numbers (see \__web_part:nn below)
% - \colon as in amsmath (MathJax's is a plain punctuation colon), and
%   \relax, which MathJax lacks and macros often end in
% - MathJax's Computer Modern font: its default, New Computer Modern, is
%   heavier than the PDF, and its Latin Modern misplaces accents on Greek
%   letters. For math the PDF sets in Times, Palatino, or New Century
%   Schoolbook, its TeX Gyre font (\web@mathfont)
% - display math that does not fit scrolls
% - math at the text size, as in TeX. Matching x-heights makes it a bit
%   small
% - for speed when MathJax typesets in the browser: math typeset as it
%   scrolls into view (ui/lazy), and no speech or braille. The semantic
%   enrichment they need took 90% of the time, 20s for a long document.
%   Readers can turn speech on in MathJax's menu. Screen readers get each
%   formula as MathML (assistiveMml, which MathJax leaves off when speech
%   is on by default)
% - speech, prerendered or from the menu, in ClearSpeak's words, with \Delta
%   as "Delta", not "triangle", and with \WebSpeechGiven, | as "given", as
%   in P(A | B). The style (\web@speech) is also set in the menu settings,
%   because the menu's setting would replace it with ClearSpeak's defaults.
%   It is expanded when the head is written, after the config
\def\web@extensions{"mathtools"}
\def\web@load{"[tex]/mathtools"}
\@for\web@package:=amscd,bbm,bboldx,braket,bussproofs,cancel,cases,centernot,%
        colortbl,dsfont,empheq,gensymb,mhchem,physics,textcomp,upgreek\do{%
    \@ifpackageloaded\web@package
        {\edef\web@extensions{\web@extensions, "\web@package"}%
         \edef\web@load{\web@load, "[tex]/\web@package"}}{}%
}
% the units extension: its \nicefrac stands for nicefrac's, and macros.lua
% defines xfrac's \sfrac with it
\let\web@units\empty
\@for\web@package:=units,nicefrac,xfrac\do{\@ifpackageloaded\web@package{\def\web@units{, "units"}}{}}
\edef\web@extensions{\web@extensions\web@units}
\edef\web@load{\web@load\ifx\web@units\empty\else, "[tex]/units"\fi}
% (newtxmath's options for other fonts turn \iftx@ntxmath off; nc is New
% Century Schoolbook)
\def\web@mathfont{tex}
\@for\web@package:=mathptmx,txfonts\do{\@ifpackageloaded\web@package{\def\web@mathfont{termes}}{}}
\@for\web@package:=mathpazo,pxfonts,newpxmath\do{\@ifpackageloaded\web@package{\def\web@mathfont{pagella}}{}}
\@ifpackageloaded{fouriernc}{\def\web@mathfont{schola}}{}
\@ifpackageloaded{newtxmath}{%
    \web@needs{iftx@ntxmath}\web@needs{iftx@nc}%
    \iftx@ntxmath\def\web@mathfont{termes}\fi
    \iftx@nc\def\web@mathfont{schola}\fi}{}
\def\web@speech{TriangleSymbol_Delta}
\newcommand\WebSpeechGiven{\gdef\web@speech{VerticalLine_Given:TriangleSymbol_Delta}}
\edef\web@mathjax{{
    "loader": {"load": ["ui/lazy", \web@load]},
    "options": {"menuOptions": {"settings": {"enrich": false, "speech": false,
        "braille": false, "assistiveMml": true,
        "speechRules": "clearspeak-\noexpand\web@speech"}},
        "sre": {"domain": "clearspeak", "style": "\noexpand\web@speech"}},
    "tex": {"tags": "none", "packages": {"[+]": [\web@extensions]},
        "macros": {"colon": "\detokenize{\\mskip2mu\\mathpunct{}\\mkern-3mu{:}\\mskip6mu}",
            "relax": ""}},
    "output": {"font": "mathjax-\web@mathfont", "displayOverflow": "scroll",
        "matchFontHeight": false}}}
% each formula at the size of the text around it, as a multiple of the body
% size (\web@bodysize): tex4ht ends the span of the text's size before each
% formula. These are html4-math.4ht's elements, with the size added.
% tex8web-text.tex turns the size off in the title, where tex4ht drops the
% fonts.
\newcommand\web@mathsize{\ifdim\f@size pt=\web@bodysize\else
    \space style="--tex8web-math-size: \web@over{\f@size pt}\web@bodysize"\fi}
\newcommand\web@mathopen[1]{%
    \ifvmode\IgnorePar\EndP\def\web@mathtag{div}\else\def\web@mathtag{span}\fi
    \edef\web@a{\noexpand\HCode{<\web@mathtag\space class="#1"\web@mathsize>}}\web@a}
\Configure{MathjaxInline}
    {\edef\web@a{\noexpand\HCode{<span class="mathjax-inline"\web@mathsize>}}\web@a}
    {\HCode{</span>}}
\Configure{MathjaxBlock}{\web@mathopen{mathjax-block}}{\HCode{</\web@mathtag>}}
\Configure{MathjaxEnv}{\web@mathopen{mathjax-env mathjax-\@currenvir}}{\HCode{</\web@mathtag>}}
% the space above and below display math, as in the PDF at the normal
% size (css/tex8web.css)
\begingroup\normalsize
    \web@Css{:root { --tex8web-display-above: \web@em\abovedisplayskip;
        --tex8web-display-below: \web@em\belowdisplayskip; }}
\endgroup
% \WebMathMacros{definitions}: definitions for MathJax that replace the
% ones macros.lua reads from the preamble and styles, for macros MathJax
% cannot read as they are. tex8web-page.tex passes them on to macros.lua
\ExplSyntaxOn
\tl_new:N \g__web_macros_tl
\NewDocumentCommand\WebMathMacros{+m}{\tl_gput_right:Nn \g__web_macros_tl {#1}}

% tex4ht escapes each formula's source (\alteqtoks, mathjax-latex-4ht.4ht)
% with four \regex_replace_all, which take most of a long document's LaTeX
% run. Here the output of \detokenize, catcode 12 characters and spaces,
% gets the same replacements, with the same tokens:
% - &, <, and > as the regex does. Its & has the catcode that & has when it
%   runs: a letter in display math, where tex4ht reads & as one
% - \par where the regex matches it: before a space, before any character
%   but a letter, a digit, or _, or at the end of the formula (its \b).
%   \detokenize ends a control word with a space, so this matches \par and
%   \\par, and leaves \partial and \parbox. A regex costs as much for a
%   formula with \partial as for any other, so each \par is found with a
%   delimited argument instead
% - \par first, because the loop reads the formula unbraced, where a
%   catcode 4 & would end a table cell
% test-mathescape.tex has every case. test/run.mjs also builds it with
% tex4ht's escaping, kept as \web@alteqtoks, and compares the pages. Not
% under fastmathjax, which escapes nothing.
\web@needs{alteqtoks}
\:CheckOption{fastmathjax}\if:Option\else
\let\web@alteqtoks\alteqtoks
\tl_const:Ne \c__web_amp_tl { \tl_to_str:n { & } }
\tl_const:Ne \c__web_lt_tl { \tl_to_str:n { < } }
\tl_const:Ne \c__web_gt_tl { \tl_to_str:n { > } }
\tl_const:Ne \c__web_par_tl { \char_generate:nn {92} {12} \tl_to_str:n { par } }
\tl_const:Ne \c__web_newlines_tl { \char_generate:nn {10} {12} \char_generate:nn {10} {12} }
\tl_const:Ne \c__web_word_tl
    { \tl_to_str:n { abcdefghijklmnopqrstuvwxyzABCDEFGHIJKLMNOPQRSTUVWXYZ0123456789_ } }
\cs_new:Npn \__web_entity:n #1 { \char_generate:nn {38} { \tex_catcode:D 38 } #1 ; }
\cs_generate_variant:Nn \tl_replace_all:Nnn { NVe }
\tl_new:N \l__web_rest_tl
\cs_new_protected:Npn \__web_par_blank: { \tl_put_right:NV \l_tmpa_tl \c__web_newlines_tl }
% \__web_par:w #1 \par #2 \q_stop, delimited by a catcode 12 \par
\cs_new_protected:Npn \__web_par_define:n #1 {
    \cs_new_protected:Npn \__web_par:w ##1 #1 ##2 \q_stop {
        \tl_put_right:Nn \l_tmpa_tl {##1}
        \tl_if_empty:nF {##2} {
            \tl_set:Nn \l__web_rest_tl {##2}
            \tl_if_head_is_space:nTF {##2} { \__web_par_blank: } {
                \exp_args:NNe \tl_if_in:NnTF \c__web_word_tl { \tl_head:n {##2} }
                    { \tl_put_right:NV \l_tmpa_tl \c__web_par_tl } { \__web_par_blank: }
            }
            \exp_after:wN \__web_par:w \l__web_rest_tl \q_stop
        }
    }
}
\exp_args:NV \__web_par_define:n \c__web_par_tl
\cs_set_protected:Npn \alteqtoks #1 {
    \tl_set:Ne \l_tmpa_tl { \detokenize {#1} }
    \tl_if_in:NVT \l_tmpa_tl \c__web_par_tl {
        % (the \par put at the end stands for the formula's end)
        \tl_set_eq:NN \l__web_rest_tl \l_tmpa_tl
        \tl_put_right:NV \l__web_rest_tl \c__web_par_tl
        \tl_clear:N \l_tmpa_tl
        \exp_after:wN \__web_par:w \l__web_rest_tl \q_stop
    }
    \tl_replace_all:NVe \l_tmpa_tl \c__web_amp_tl { \__web_entity:n { amp } }
    \tl_replace_all:NVe \l_tmpa_tl \c__web_lt_tl { \__web_entity:n { lt } }
    \tl_replace_all:NVe \l_tmpa_tl \c__web_gt_tl { \__web_entity:n { gt } }
    \tl_set:Ne \l_tmpb_tl { \l_tmpa_tl }
    \HCode { \l_tmpb_tl }
}
\fi

% siunitx. tex4ht writes each number and unit siunitx prints as a separate
% formula (\__siunitx_print_math_auxii:n, html4.4ht). A quantity in text
% is \ensuremath around an internal command, which tex4ht's mathjax option
% passes to MathJax unchanged, and MathJax knows none of siunitx's commands
% in a formula. So each quantity in text, and each siunitx command in a
% formula, is run in a box, and what it prints is kept as math: numbers,
% units, and product. A quantity in text is then one formula, as in the
% PDF. A quantity siunitx sets as text is left to tex4ht. Ranges, lists,
% and products in text are left to tex4ht too, as numbers in a sentence.
% The unit \percent prints a %, which starts a comment for MathJax.
\@ifpackageloaded{siunitx}{
    \web@needs{l__siunitx_quantity_product_tl}
    \web@needs{siunitx_quantity_print:nn}
    \web@needs{l_siunitx_number_mode_str}
    \web@needs{l_siunitx_unit_mode_str}
    \tl_new:N \l__web_si_tl
    \tl_new:N \g__web_si_tl
    \str_new:N \l__web_si_str
    \bool_new:N \l__web_si_bool
    \box_new:N \l__web_si_box
    \quark_new:N \q__web_si
    \clist_new:N \l__web_si_clist
    % (after tex4ht's definition, made in begindocument under \fixmathjaxtoc)
    \web@needs{fixmathjaxtoc}
    \web@needs{__siunitx_print_math_auxii:n}
    \cs_generate_variant:Nn \str_replace_all:Nnn { Nee }
    % - every % and \% as \%, in a string, because \tl_replace_all:Nnn
    %   skips text in braces
    % - a \, between digits, siunitx's digit group separator, as an empty mn
    %   of its width: MathJax's speech reads an mspace between two numbers
    %   as "times"
    % - \micro's \text{\textmu} (pdfTeX) as \text{μ}: MathJax shows
    %   \textmu as it is, and sets a micro sign (U+00B5) italic and a Greek
    %   mu upright, as in the PDF
    \regex_const:Nn \c__web_si_digits_regex { ([0-9]) \\, ([0-9]) }
    \str_const:Ne \c__web_si_textmu_str { \tl_to_str:n { \text { \textmu } } }
    \str_const:Ne \c__web_si_mu_str { \tl_to_str:n { \text { μ } } }
    \str_const:Ne \c__web_si_thin_str { \tl_to_str:n { \, } }
    % siunitx's own math, which tex4ht replaces with source for MathJax
    % (html4.4ht), for the glyphs of a saved box below. As in siunitx,
    % without its cases for \sfrac and for siunitx's own HTML
    \web@needs{__siunitx_print_math_text_setup:}
    \web@needs{__siunitx_print_math_auxiii:n}
    \cs_new_protected:Npn \__web_saved_si_math:n #1 {
        \group_begin:
            \__siunitx_print_math_text_setup:
            \__siunitx_print_math_auxiii:n {#1}
        \group_end:
    }
    \AddToHook { begindocument/end } {
        \cs_set_eq:NN \__web_si_auxii:n \__siunitx_print_math_auxii:n
        \cs_gset_protected:Npn \__siunitx_print_math_auxii:n #1 {
            \str_set:Nn \l_tmpa_str {#1}
            \str_replace_all:Nee \l_tmpa_str { \c_backslash_str \c_percent_str } { \c_percent_str }
            \exp_args:NNV \str_replace_all:Nnn \l_tmpa_str \c_percent_str { \% }
            \exp_args:NNV \str_if_in:NnT \l_tmpa_str \c__web_si_thin_str {
                \regex_replace_all:NnN \c__web_si_digits_regex
                    { \1 \\mmlToken\{mn\}[style="padding-left:0.167em"]\{\} \2 } \l_tmpa_str
            }
            \exp_args:NNVV \str_replace_all:Nnn \l_tmpa_str \c__web_si_textmu_str \c__web_si_mu_str
            \a:MathjaxInline \exp_args:NV \__web_si_auxii:n \l_tmpa_str \b:MathjaxInline
        }
    }
    % \g__web_si_tl: what #1 prints
    \cs_new_protected:Npn \__web_si_keep:n #1 { \tl_gput_right:Nn \g__web_si_tl {#1} }
    \cs_new_protected:Npn \__web_si_print:n #1 {
        \tl_gclear:N \g__web_si_tl
        \hbox_set:Nn \l__web_si_box {
            \cs_set_protected:Npn \__web_si_auxii:n ##1 { \__web_si_keep:n {##1} }
            \cs_set_eq:NN \siunitx_quantity_print:nn \__web_si_quantity:nn
            \cs_set_eq:NN \ensuremath \use:n
            \keys_set:ne { siunitx } {
                mode = math ,
                quantity-product = { \__web_si_keep:n { \exp_not:V \l__siunitx_quantity_product_tl } }
            }
            \cs_set_protected:Npn \__siunitx_compound_print_separator:n ##1 {
                \bool_if:NTF \l__siunitx_compound_separator_text_bool
                    { \__web_si_text:n {##1} } { \__web_si_separator:n {##1} }
            }
            \cs_set_protected:Npn \__siunitx_compound_print_boundary:n ##1 {
                \bool_if:NTF \l__siunitx_compound_boundary_text_bool
                    { \__web_si_text:n {##1} } { \__web_si_boundary:n {##1} }
            }
            #1
        }
    }
    % the text of ranges, lists, and products (their phrases and separators,
    % and a range's opening phrase), kept as one \text with the spaces
    % around it, which MathJax keeps there and leaves out of the speech.
    % (MathJax reads a \ as "empty", and tex4ht's \ is a box, so a \ in
    % a phrase is a space.) Translations are expanded, font commands
    % dropped, and dashes made characters, as MathJax sets -- in \text as
    % two hyphens. Any other text command is left as it is, which MathJax
    % shows in red. A phrase that is empty once expanded, such as the
    % default range-open-phrase, is left out
    \web@needs{__siunitx_compound_print_separator:n}
    \web@needs{l__siunitx_compound_separator_text_bool}
    \web@needs{l__siunitx_compound_boundary_text_bool}
    \web@needs{__siunitx_compound_print_boundary:n}
    \cs_set_eq:NN \__web_si_separator:n \__siunitx_compound_print_separator:n
    \cs_set_eq:NN \__web_si_boundary:n \__siunitx_compound_print_boundary:n
    \cs_new_protected:Npn \__web_si_text:n #1 {
        \group_begin:
            \cs_set_eq:NN \TextOrMath \use_i:nn
            \cs_set_eq:NN \text \use:n
            \cs_set:Npn \ { ~ }
            \text_declare_purify_equivalent:Nn \textendash { – }
            \text_declare_purify_equivalent:Nn \textemdash { — }
            \tl_gset:Ne \g__web_si_phrase_tl { \text_purify:n {#1} }
        \group_end:
        \tl_greplace_all:Nnn \g__web_si_phrase_tl { --- } { — }
        \tl_greplace_all:Nnn \g__web_si_phrase_tl { -- } { – }
        \tl_if_blank:VF \g__web_si_phrase_tl {
            \__web_si_keep:e { \exp_not:N \text { \exp_not:V \g__web_si_phrase_tl } }
        }
    }
    \tl_new:N \g__web_si_phrase_tl
    \cs_generate_variant:Nn \__web_si_keep:n { e }
    \cs_set_eq:NN \__web_si_quantity:nn \siunitx_quantity_print:nn
    % siunitx puts this protected function, unexpanded, in the math of a
    % complex quantity in brackets, where it runs when typeset. MathJax
    % gets the brackets instead.
    \web@needs{__siunitx_complex_format_bracket:n}
    \cs_set:Npn \__siunitx_complex_format_bracket:n #1 {
        \exp_not:V \l__siunitx_complex_bracket_open_tl
        #1
        \exp_not:V \l__siunitx_complex_bracket_close_tl
    }
    % (siunitx's test for whether it uses \ensuremath)
    \cs_set_protected:Npn \siunitx_quantity_print:nn #1#2 {
        \bool_lazy_any:nTF {
            { \mode_if_math_p: }
            { ! \str_if_eq_p:Vn \l_siunitx_number_mode_str { math } }
            { ! \str_if_eq_p:Vn \l_siunitx_unit_mode_str { math } }
        } { \__web_si_quantity:nn {#1} {#2} } {
            \__web_si_print:n { \__web_si_quantity:nn {#1} {#2} }
            \__web_saved:n { \__web_si_quantity:nn {#1} {#2} }
            \a:MathjaxInline
            \use:e { \alteqtoks { \exp_not:N \( \exp_not:V \g__web_si_tl \exp_not:N \) } }
            \b:MathjaxInline
        }
    }
    \cs_set_eq:NN \__web_si_alteqtoks:n \alteqtoks
    \cs_set_protected:Npn \alteqtoks #1 {
        \str_set:Nn \l__web_si_str {#1}
        \bool_set_false:N \l__web_si_bool
        \clist_map_inline:Nn \l__web_si_clist {
            \str_if_in:NnT \l__web_si_str {##1}
                { \bool_set_true:N \l__web_si_bool \clist_map_break: }
        }
        \bool_if:NTF \l__web_si_bool
            {
                % (a display comes as its source, detokenized. & is
                % rescanned as a letter, so the walk can read it unbraced)
                \tl_set_rescan:Nno \l_tmpa_tl { \char_set_catcode_letter:N \& } \l__web_si_str
                \exp_args:NV \__web_si_walk:n \l_tmpa_tl
                \exp_args:NV \__web_si_alteqtoks:n \l__web_si_tl
            }
            { \__web_si_alteqtoks:n {#1} }
    }
    % \l__web_si_tl: #1, with siunitx's commands replaced by what they print
    \cs_new_protected:Npn \__web_si_walk:n #1 {
        \tl_clear:N \l__web_si_tl
        \__web_si_loop: #1 \q__web_si
    }
    \cs_new_protected:Npn \__web_si_loop: {
        \peek_N_type:TF { \__web_si_token:N }
            { \peek_catcode:NTF \c_group_begin_token { \__web_si_group:n } { \__web_si_space:w } }
    }
    \cs_new_protected:Npn \__web_si_token:N #1 {
        \token_if_eq_meaning:NNF #1 \q__web_si {
            \cs_if_exist:cTF { __web_si_ \token_to_str:N #1 }
                { \use:c { __web_si_ \token_to_str:N #1 } }
                { \tl_put_right:Nn \l__web_si_tl {#1} \__web_si_loop: }
        }
    }
    \cs_new_protected:Npn \__web_si_group:n #1 {
        \group_begin:
            \__web_si_walk:n {#1}
            \use:e {
                \group_end:
                \tl_put_right:Nn \exp_not:N \l__web_si_tl { { \exp_not:V \l__web_si_tl } }
            }
        \__web_si_loop:
    }
    \use:n { \cs_new_protected:Npn \__web_si_space:w } ~ {
        \tl_put_right:Nn \l__web_si_tl { ~ }
        \__web_si_loop:
    }
    % each command, with its arguments as siunitx reads them. Not \qty under
    % physics, because siunitx then leaves \qty to physics
    \cs_new_protected:Npn \__web_si_command:Nnn #1#2#3 {
        \clist_put_right:Nn \l__web_si_clist {#1}
        \exp_args:Nc \NewDocumentCommand { __web_si_ \token_to_str:N #1 } {#2} {
            \use:e { \__web_si_print:n { \exp_not:N #1 #3 } }
            \tl_put_right:NV \l__web_si_tl \g__web_si_tl
            \__web_si_loop:
        }
    }
    \clist_map_inline:nn { \num, \unit, \si, \ang, \numlist, \numproduct, \complexnum }
        { \__web_si_command:Nnn #1 { O{} m } { [ \exp_not:n {##1} ] { \exp_not:n {##2} } } }
    \clist_map_inline:nn { \numrange, \qtylist, \SIlist, \qtyproduct, \complexqty }
        {
            \__web_si_command:Nnn #1 { O{} m m }
                { [ \exp_not:n {##1} ] { \exp_not:n {##2} } { \exp_not:n {##3} } }
        }
    \@ifpackageloaded{physics}{}{
        \__web_si_command:Nnn \qty { O{} m m }
            { [ \exp_not:n {#1} ] { \exp_not:n {#2} } { \exp_not:n {#3} } }
    }
    \clist_map_inline:nn { \qtyrange, \SIrange }
        {
            \__web_si_command:Nnn #1 { O{} m m m }
                { [ \exp_not:n {##1} ] { \exp_not:n {##2} } { \exp_not:n {##3} } { \exp_not:n {##4} } }
        }
    \__web_si_command:Nnn \SI { O{} m o m }
        { [ \exp_not:n {#1} ] { \exp_not:n {#2} } \IfValueT {#3} { [ \exp_not:n {#3} ] } { \exp_not:n {#4} } }
}{}

% mhchem's \ce, as an inline formula. tex4ht writes its source as text in
% \texttt (mathjax-latex-4ht.4ht), and ends and starts the font's span at
% each line break of the HTML. MathJax finds no formula whose \) is in
% another element, and make4ht's joincharacters filter joins the spans
% only after prerendering. In a heading, \ce is written to the .aux and
% the table of contents as it stands.
\@ifpackageloaded{mhchem}{
    \cs_new_protected:Npn \__web_ce:n #1 { \( \ce {#1} \) }
    \cs_new:Npn \__web_ce: {
        \token_if_eq_meaning:NNTF \protect \@typeset@protect { \__web_ce:n } { \exp_not:N \ce }
    }
    \AddToHook { begindocument/end } { \HLet \ce \__web_ce: }
}{}

% a picture in display math is typeset by TeX, rescanning the source,
% which tex4ht read with & as a letter: a body that is only one
% tikzpicture, \tikz, tikzcd, \includegraphics, or picture environment,
% maybe in a \resizebox or a \scalebox, with \label, \notag, and
% \nonumber anywhere around it, \qedhere and \tag after it, and a period
% or a comma after it, maybe after one spacing command (\, \: \; \!
% \quad \qquad or a control space); and a \[ \] with a tikzcd. MathJax knows none of
% them. Other pictures beside math, and a body with two pictures, stay
% MathJax's. amsmath's displays take a \tag, and so does \[ \] with
% amsmath; $$ $$, displaymath, and eqnarray do not.
\tl_const:Nn \c__web_picture_tl {
    (?: \\begin \s* \{ tikzpicture \} .* \\end \s* \{ tikzpicture \}
        | \\begin \s* \{ tikzcd \} .* \\end \s* \{ tikzcd \}
        | \\begin \s* \{ picture \} .* \\end \s* \{ picture \}
        | \\includegraphics \s* \*? \s* (?: \[ [^\]]* \] \s* )* \{ [^\{\}]* \}
        | \\tikz \b \s* (?: \[ [^\]]* \] \s* )?
            (?: \{ (?: [^\{\}] | \{ (?: [^\{\}] | \{ [^\{\}]* \} )* \} )* \} | (?: [^;\{\}] | \{ (?: [^\{\}] | \{ (?: [^\{\}] | \{ [^\{\}]* \} )* \} )* \} )* ; ) )
}
\tl_const:Nn \c__web_marks_tl {
    \s | \\label \s* \{ [^\{\}]* \} | \\notag \b | \\nonumber \b
}
\exp_args:NNe \regex_const:Nn \c__web_picture_regex {
    \exp_not:n { \A (?: } \exp_not:V \c__web_marks_tl \exp_not:n { )* (?: }
        \exp_not:V \c__web_picture_tl
        \exp_not:n { | (?: \\resizebox \s* \*? \s* \{ [^\{\}]* \} \s* \{ [^\{\}]* \}
            | \\scalebox \s* \{ [^\{\}]* \} (?: \s* \[ [^\]]* \] )? ) \s* \{ \s* }
        \exp_not:V \c__web_picture_tl
    \exp_not:n { \s* \} ) (?: } \exp_not:V \c__web_marks_tl
    \exp_not:n { | \\qedhere \b | \\tag \s* \*? \s* \{ [^\{\}]* \} )*
        (?: (?: (?: \\[,:;!\ ] | \\q?quad \b ) \s* )? [.,] (?: }
    \exp_not:V \c__web_marks_tl
    \exp_not:n { | \\qedhere \b | \\tag \s* \*? \s* \{ [^\{\}]* \} )* )? \Z }
}
% a period or comma after the picture, with its spacing command before
% it and maybe marks after it. They are taken out of the body, and \__web_picture_punct: sets it after the
% picture, on its bottom line (css/tex8web.css). (The body is read again
% as it is set, which a tikzcd's & needs)
\regex_const:Nn \c__web_picture_punct_regex {
    ( (?: \c{ (?: [,:;!\ ] | q?quad ) } \s* )? [.,] ) ( (?: \s | \c{label} \s* \cB. (?: \c[^BE]. )* \cE. | \c{notag} | \c{nonumber} | \c{qedhere}
        | \c{tag} \s* (?: \* \s* )? \cB. (?: \c[^BE]. )* \cE. )* ) \Z
}
\tl_new:N \l__web_picture_body_tl
\tl_new:N \l__web_picture_punct_tl
\seq_new:N \l__web_picture_seq
\cs_new_protected:Npn \__web_picture_body:n #1 {
    \tl_set:Nn \l__web_picture_body_tl {#1}
    \tl_clear:N \l__web_picture_punct_tl
    \regex_extract_once:NnNT \c__web_picture_punct_regex {#1} \l__web_picture_seq {
        \tl_set:Ne \l__web_picture_punct_tl { \seq_item:Nn \l__web_picture_seq { 2 } }
        \regex_replace_once:NnN \c__web_picture_punct_regex { \2 } \l__web_picture_body_tl
    }
}
\cs_new_protected:Npn \__web_picture_punct: {
    \tl_if_empty:NF \l__web_picture_punct_tl {
        \HCode{<span~class="tex8web-picture-punct">}
        \hbox { \tl_use:N \l__web_picture_punct_tl }
        \HCode{</span>}
    }
}
% the end of each picture, to find a body with two
\regex_const:Nn \c__web_picture_end_regex
    { \\end \s* \{ (?: tikzpicture | tikzcd | picture ) \} | \\tikz \b }
\str_new:N \l__web_picture_str
\int_new:N \l__web_picture_int
% (the regexes run only on a body with a picture's command. \str_if_in
% is needed for the \begin{...} commands: \tl_if_in is true for any body
% when the text it looks for has braces)
\prg_new_protected_conditional:Npnn \__web_if_picture:n #1 { TF } {
    \str_set:Nn \l__web_picture_str {#1}
    \__web_if_picture_command:TF
        {
            \exp_args:NNV \regex_count:NnN \c__web_picture_end_regex \l__web_picture_str \l__web_picture_int
            \int_compare:nNnTF \l__web_picture_int < 2
                {
                    \exp_args:NNV \regex_match:NnTF \c__web_picture_regex \l__web_picture_str
                        { \prg_return_true: } { \prg_return_false: }
                }
                { \prg_return_false: }
        }
        { \prg_return_false: }
}
\cs_new_protected:Npn \__web_if_picture_command:TF #1#2 {
    \str_if_in:NnTF \l__web_picture_str {\includegraphics} {#1} {
        \str_if_in:NnTF \l__web_picture_str {\tikz} {#1} {
            \str_if_in:NnTF \l__web_picture_str {\begin{tikzpicture}} {#1} {
                \str_if_in:NnTF \l__web_picture_str {\begin{tikzcd}} {#1}
                    { \str_if_in:NnTF \l__web_picture_str {\begin{picture}} {#1} {#2} }
            }
        }
    }
}
% a picture with a \tag in a \[ \], which takes one with amsmath, as
% equation*
\prg_new_protected_conditional:Npnn \__web_if_tagged_picture:n #1 { TF } {
    \cs_if_exist:NTF \tag
        {
            \tl_if_in:nnTF {#1} {\tag}
                { \__web_if_picture:nTF {#1} { \prg_return_true: } { \prg_return_false: } }
                { \prg_return_false: }
        }
        { \prg_return_false: }
}
% the same, in a display that takes no \tag: \[ \], $$ $$, and displaymath
\prg_new_protected_conditional:Npnn \__web_if_bare_picture:n #1 { TF } {
    \__web_if_picture:nTF {#1}
        { \tl_if_in:nnTF {#1} {\tag} { \prg_return_false: } { \prg_return_true: } }
        { \prg_return_false: }
}
\cs_new_protected:Npn \__web_picture:n #1 {
    \char_set_catcode_alignment:N \& \tex_scantokens:D {#1}
}
% A \[ \] with a \qedhere is walked for its box (see \__web_part:nn)
\web@needs{AltlDisplay}
\web@needs{AltlDisplayDollars}
\web@needs{VerbMathToks}
\cs_set_eq:NN \__web_AltlDisplay:w \AltlDisplay
\cs_set_eq:NN \__web_AltlDisplayDollars:w \AltlDisplayDollars
% (set as in a display environment below, with no number: the picture in
% an \hbox, then a period or comma after it, then the box of a \qedhere.
% tex4ht's mathjax option leaves TeX in text mode at \[ and $$, so after
% a blank line an unboxed picture would start a paragraph in the display)
\cs_new_protected:Npn \__web_picture_bare:n #1 {
    \group_begin:
        \tl_clear:N \l__web_picture_tag_tl
        \bool_set_false:N \l__web_picture_back_bool
        \__web_picture_set:n {#1}
    \group_end:
}
\cs_set:Npn \AltlDisplayDollars #1 $$ {
    \__web_if_bare_picture:nTF {#1}
        { \__web_picture_bare:n {#1} $$ }
        { \__web_saved:n {#1} \__web_AltlDisplayDollars:w #1 $$ }
}
% (a tikzcd beside other math is set in the display's math)
\cs_set:Npn \AltlDisplay #1 \] {
    \__web_if_bare_picture:nTF {#1}
        { \__web_picture_bare:n {#1} \] }
        {
            \__web_if_tagged_picture:nTF {#1}
                { \__web_picture_tagged:n {#1} \] }
                {
                    \str_if_in:nnTF {#1} {\begin{tikzcd}}
                        { \__web_picture:n {#1\]} }
                        {
                            \str_if_in:nnTF {#1} {\qedhere}
                                {
                                    \__web_saved_qed:n {#1}
                                    \tl_clear:N \l__web_alignment_tl
                                    \__web_rows:nn {equation*} {#1}
                                    \exp_after:wN \__web_AltlDisplay:w \l__web_body_tl \]
                                }
                                { \__web_saved:n { \[#1\] } \__web_AltlDisplay:w #1 \] }
                        }
                }
        }
}

% amsmath's displays and LaTeX's eqnarray, which tex4ht passes to MathJax
% as they are, except:
% - MathJax has no \intertext, so the alignment is split there. MathJax
%   gets the parts, and the text between them is typeset here, in a group,
%   because amsmath sets it in a \vbox that ends its font commands. Each
%   part starts with \webalignment{n}, n numbering the split alignments,
%   and tex8web.js gives the parts of one alignment the same columns. Each
%   part of an alignat gets its number of columns
% - equation numbers: MathJax would number the equations from 1, while
%   \eqref and \ref print TeX's numbers, and tex4ht steps the counter once
%   per amsmath environment, and never in an eqnarray. So each row is
%   numbered here, and MathJax gets the number as a \tag
% - amsmath starts each right-hand column (&= b, &-c) with {}, so a
%   relation or binary operator there is spaced as one. MathJax does not
\int_new:N \g__web_alignment_int
\tl_new:N \l__web_alignment_tl
\cs_new_protected:Npn \__web_display:nn #1#2 {
    \tl_clear:N \l__web_alignment_tl
    \tl_if_in:nnTF {#2} {\intertext}
        {
            \int_gincr:N \g__web_alignment_int
            \tl_set:Ne \l__web_alignment_tl
                { \exp_not:N \webalignment { \int_use:N \g__web_alignment_int } }
            \str_if_in:nnTF {#1} {alignat}
                { \__web_alignat:nw {#1} #2 \q_stop }
                { \__web_alignment:nn {#1} {#2} }
        }
        { \__web_alignment:nn {#1} {#2} }
}
\cs_new_protected:Npn \__web_alignat:nw #1#2#3 \q_stop {
    \tl_put_right:Nn \l__web_alignment_tl { {#2} }
    \__web_alignment:nn {#1} {#3}
}
\cs_new_protected:Npn \__web_alignment:nn #1#2 {
    \tl_if_in:nnTF {#2} {\intertext}
        { \__web_intertext:nw {#1} #2 \q_stop }
        { \__web_part:nn {#1} {#2} }
}
\tl_new:N \l__web_before_tl
\cs_new_protected:Npn \__web_intertext:nw #1#2 \intertext #3#4 \q_stop {
    \tl_set:Nn \l__web_before_tl {#2}
    \regex_replace_once:nnN { \c{\\} \s* \Z } {} \l__web_before_tl
    \tl_if_blank:VF \l__web_before_tl { \exp_args:Nno \__web_part:nn {#1} {\l__web_before_tl} }
    \HCode{<span~class="tex8web-intertext">} \group_begin: #3 \group_end: \HCode{</span>}
    \tl_if_blank:nF {#4} { \__web_alignment:nn {#1} {#4} }
}

% a part of a display: its tokens are walked to number the rows, and the
% part is passed to MathJax as tex4ht's mathjax option does. A \[ \] with a
% \qedhere is walked as equation*, but stays a \[ \] for MathJax
\tl_new:N \l__web_body_tl
\tl_new:N \l__web_label_tl
\tl_new:N \l__web_tag_tl
\bool_new:N \l__web_numbered_bool
\bool_new:N \l__web_rows_bool
\bool_new:N \l__web_number_row_bool
\bool_new:N \l__web_own_tag_bool
\bool_new:N \l__web_qed_wanted_bool
\bool_new:N \l__web_qed_bool
\tl_new:N \l__web_qed_tag_tl
\tl_new:N \l__web_qed_symbol_tl
\seq_new:N \l__web_envs_seq
\seq_new:N \l__web_columns_seq
\tl_new:N \l__web_next_tl
\tl_new:N \l__web_count_tl
\tl_new:N \l__web_env_tl
% (tex8web.js has a table of these environments too, for their row spacing)
\clist_const:Nn \c__web_numbered_clist { equation, multline, align, gather, alignat, flalign, eqnarray }
\clist_const:Nn \c__web_rows_clist { align, gather, alignat, flalign, eqnarray, eqnarray* }
\clist_const:Nn \c__web_alignments_clist
    { align, align*, alignat, alignat*, flalign, flalign*, split, aligned, alignedat }
\cs_new_protected:Npn \__web_part:nn #1#2 {
    \__web_rows:nn {#1} {#2}
    \exp_args:No \__web_mathjax:nn {\l__web_body_tl} {#1}
}
\cs_new_protected:Npn \__web_rows:nn #1#2 {
    \tl_set_eq:NN \l__web_body_tl \l__web_alignment_tl
    \seq_clear:N \l__web_envs_seq
    \seq_clear:N \l__web_columns_seq
    \seq_push:Nn \l__web_envs_seq {#1}
    \seq_push:Nn \l__web_columns_seq {0}
    \clist_if_in:NnTF \c__web_numbered_clist {#1}
        { \bool_set_true:N \l__web_numbered_bool }
        { \bool_set_false:N \l__web_numbered_bool }
    \clist_if_in:NnTF \c__web_rows_clist {#1}
        { \bool_set_true:N \l__web_rows_bool }
        { \bool_set_false:N \l__web_rows_bool }
    \__web_row_start:
    \tl_clear:N \l__web_next_tl
    \__web_walk:w #2 \q_nil
    \__web_row_end:
}
\cs_new_protected:Npn \__web_mathjax:nn #1#2 {
    \NoFonts\expandafter\VerbMathToks\expandafter{\detokenize{#1}}{#2}\EndNoFonts
}
\cs_new_protected:Npn \__web_row_start: {
    \bool_set_true:N \l__web_number_row_bool
    \bool_set_false:N \l__web_own_tag_bool
    \bool_set_false:N \l__web_qed_wanted_bool
    \bool_set_false:N \l__web_qed_bool
    \tl_clear:N \l__web_label_tl
    \tl_clear:N \l__web_tag_tl
}
% the number at the end of a row, or of a display numbered as a whole
\cs_new_protected:Npn \__web_row_end: {
    \tl_if_empty:NF \l__web_tag_tl
        { \protected@edef\@currentlabel{\l__web_tag_tl} }
    \bool_lazy_and:nnT \l__web_qed_wanted_bool { \bool_not_p:n \l__web_own_tag_bool }
        {
            \__web_qedhere:
            \bool_lazy_and:nnT \l__web_qed_bool \l__web_rows_bool
                { \bool_set_false:N \l__web_number_row_bool }
        }
    \bool_lazy_and:nnTF \l__web_numbered_bool \l__web_number_row_bool
        {
            \refstepcounter{equation}
            \bool_if:NTF \l__web_qed_bool
                { \__web_qed_row:e { (\theequation) \exp_not:N \\ } }
                {
                    \tl_put_right:Ne \l__web_body_tl { \exp_not:N \tag {\theequation} }
                    \__web_tag_label:
                }
        }
        {
            \bool_if:NTF \l__web_qed_bool
                { \__web_qed_row:n {} }
                { \bool_if:NT \l__web_own_tag_bool { \__web_tag_label: } }
            \__web_eqnarray_label:
        }
    \tl_if_empty:NF \l__web_label_tl { \exp_args:No \label {\l__web_label_tl} }
    \__web_row_start:
}
% MathJax gets a row's \label only after the row's tag, whose element the
% \label names. MathJax keeps a \label until a tag takes it (makeTag,
% Tags.ts), so a \label on an untagged row would be given to a later tag,
% or give "Multiple \label" with the next \label
\cs_new_protected:Npn \__web_tag_label: {
    \tl_if_empty:NF \l__web_label_tl {
        \tl_put_right:Ne \l__web_body_tl
            { \exp_not:N \label { \exp_not:V \l__web_label_tl } }
    }
}
% A \label on an unnumbered row of an eqnarray, or in an eqnarray*, gets the
% number the next row would get, as in the PDF, where the counter is stepped
% after each numbered row (\@@eqncr, latex.ltx). It leads to an anchor of
% its own, which \refstepcounter makes
\prg_new_conditional:Npnn \__web_eqnarray: { p } {
    \str_case_e:nnTF { \seq_item:Nn \l__web_envs_seq { -1 } }
        { {eqnarray} {} {eqnarray*} {} } \prg_return_true: \prg_return_false:
}
\cs_new_protected:Npn \__web_eqnarray_label: {
    \bool_lazy_and:nnT
        { \__web_eqnarray_p: }
        { ! \tl_if_empty_p:N \l__web_label_tl }
        {
            \group_begin:
                \refstepcounter{equation}
                \exp_args:No \label {\l__web_label_tl}
                \addtocounter{equation}{-1}
            \group_end:
            \tl_clear:N \l__web_label_tl
        }
}
\cs_new_protected:Npn \__web_walk:w {
    \peek_remove_spaces:n {
        \peek_catcode:NTF \c_group_begin_token
            { \__web_group:n } { \__web_token:N }
    }
}
\cs_new_protected:Npn \__web_group:n #1 {
    \str_if_eq:VnTF \l__web_next_tl {label}
        { \tl_set:Nn \l__web_label_tl {#1} }
        {
            \str_case:Vn \l__web_next_tl {
                {begin} {
                    \seq_push:Nn \l__web_envs_seq {#1}
                    \seq_push:Nn \l__web_columns_seq {0}
                }
                {tag} { \tl_set:Nn \l__web_tag_tl {#1} }
            }
            \tl_put_right:Nn \l__web_body_tl { {#1} }
        }
    \tl_clear:N \l__web_next_tl
    \__web_walk:w
}
\cs_new_protected:Npn \__web_token:N #1 {
    \token_if_eq_meaning:NNF #1 \q_nil {
        \tl_clear:N \l__web_next_tl
        \str_case_e:nnF { \token_to_str:N #1 }
            { { \token_to_str:N \\ } {} { \token_to_str:N \label } {} }
            { \tl_put_right:Nn \l__web_body_tl {#1} }
        \token_if_eq_meaning:NNTF #1 & { \__web_column: } {
            \str_case_e:nn { \token_to_str:N #1 } {
                { \token_to_str:N \begin } { \tl_set:Nn \l__web_next_tl {begin} }
                { \token_to_str:N \end } {
                    \seq_pop:NN \l__web_envs_seq \l__web_env_tl
                    \seq_pop:NN \l__web_columns_seq \l__web_count_tl
                }
                { \token_to_str:N \label } { \tl_set:Nn \l__web_next_tl {label} }
                { \token_to_str:N \tag } {
                    \bool_set_false:N \l__web_number_row_bool
                    \bool_set_true:N \l__web_own_tag_bool
                    \tl_set:Nn \l__web_next_tl {tag}
                }
                { \token_to_str:N \notag } { \bool_set_false:N \l__web_number_row_bool }
                { \token_to_str:N \nonumber } { \bool_set_false:N \l__web_number_row_bool }
                { \token_to_str:N \qedhere } { \__web_qed_wanted: }
                { \token_to_str:N \\ } { \__web_newline: }
            }
        }
        \__web_walk:w
    }
}
% \\: a new row, with a number in the environments numbered by row
\cs_new_protected:Npn \__web_newline: {
    \seq_pop:NN \l__web_columns_seq \l__web_count_tl
    \seq_push:Nn \l__web_columns_seq {0}
    \bool_lazy_and:nnT \l__web_rows_bool
        { \int_compare_p:nNn { \seq_count:N \l__web_envs_seq } = 1 }
        { \__web_row_end: }
    \tl_put_right:Nn \l__web_body_tl { \\ }
}
% &: in an alignment, {} at the start of each right-hand column
\cs_new_protected:Npn \__web_column: {
    \seq_pop:NN \l__web_columns_seq \l__web_count_tl
    \seq_push:Ne \l__web_columns_seq { \int_eval:n { \l__web_count_tl + 1 } }
    \seq_get:NN \l__web_envs_seq \l__web_env_tl
    \int_if_odd:nT { \l__web_count_tl + 1 } {
        \clist_if_in:NVT \c__web_alignments_clist \l__web_env_tl
            { \tl_put_right:Nn \l__web_body_tl { {} } }
    }
}
% amsthm's \qedhere in a display, which TeX never runs, because MathJax gets
% the source. Without this code, the proof's box is set below the display.
% At the end of the row, or of the \[ \] display, the box is taken off the
% proof's stack, as \qedhere does, and the display gets the box as a tag
% where amsthm puts it: at the right of the last row, which then has no
% number, or under an equation's number. In a row with its own \tag, the
% box stays on the stack and is set below the display. The tag gets the
% row's \label, or one of its own, because MathJax names an unlabelled
% tag's element after its text, and two boxes would share an id. The box
% is the \qedsymbol in force, chosen as for the text (\web@qed@choose,
% tex8web-text.tex)
\int_new:N \g__web_qed_int
% In an eqnarray, amsthm sets the box in the row, which keeps any number
% (\math@qedhere, amsthm.sty). Here the box stays on the stack, and is set
% below the display
\cs_new_protected:Npn \__web_qed_wanted: {
    \bool_if:nF { \__web_eqnarray_p: } { \bool_set_true:N \l__web_qed_wanted_bool }
}
\@ifpackageloaded{amsthm}{ \web@needs{QED@stack} }{}
\cs_new_protected:Npn \__web_qedhere: {
    \cs_if_exist:NT \QED@stack {
        \tl_if_empty:NF \QED@stack
            { \exp_last_unbraced:NV \__web_qedhere:Nnw \QED@stack \q_stop }
    }
}
\cs_new_protected:Npn \__web_qedhere:Nnw #1#2#3 \q_stop {
    \tl_if_empty:nF {#2} {
        \bool_set_true:N \l__web_qed_bool
        \tl_gset:Nn \QED@stack { #1 {} #3 }
    }
}
\cs_new_protected:Npn \__web_qed_tag:n #1 {
    \web@qed@choose
    \int_gincr:N \g__web_qed_int
    \tl_set:Ne \l__web_qed_symbol_tl {
        \ifx \web@qedsym \web@openbox $ \exp_not:N \square $
        \else \exp_not:o { \web@qedsym } \fi
    }
    \tl_set:Ne \l__web_qed_tag_tl {
        \exp_not:N \tag * { \exp_not:n {#1} \exp_not:V \l__web_qed_symbol_tl }
        \exp_not:N \label {
            \tl_if_empty:NTF \l__web_label_tl
                { web-qed- \int_use:N \g__web_qed_int } { \exp_not:V \l__web_label_tl }
        }
    }
}
\cs_new_protected:Npn \__web_qed_row:n #1 {
    \__web_qed_tag:n {#1}
    \tl_put_right:NV \l__web_body_tl \l__web_qed_tag_tl
}
\cs_generate_variant:Nn \__web_qed_row:n { e }
% a display environment whose body is only a picture is set as in a \[ \]
% (see \c__web_picture_regex), in an \hbox, so that after a blank line it
% starts no paragraph, as MathJax's display does not. Its number, or its
% \tag, follows the picture in the display, and then the box of a
% \qedhere, as amsthm sets the box under the number (css/tex8web.css)
\tl_new:N \l__web_picture_tag_tl
\bool_new:N \g__web_picture_qed_bool
% (the number's step comes before the display, so that its anchor does,
% as before MathJax's displays: page.lua marks it, and a link to it shows
% the whole display)
\cs_new_protected:Npn \__web_picture_display:nn #1#2 {
    \group_begin:
        \__web_picture_number:nn {#1} {#2}
        \a:MathjaxBlock
        \__web_picture_set:n {#2}
        \b:MathjaxBlock
    \group_end:
}
% a \[ \] with a \tag, in its math: tex4ht has opened the display, which is
% closed before the tag's anchor and opened again after it, as the same
% element. page.lua removes the empty display this leaves
\cs_new_protected:Npn \__web_picture_tagged:n #1 {
    \group_begin:
        \b:MathjaxBlock
        \__web_picture_number:nn {equation*} {#1}
        \a:MathjaxBlock
        \__web_picture_set:n {#1}
    \group_end:
}
% the display's number or tag, in \l__web_picture_tag_tl, with the step of
% the counter that a \label names
\cs_new_protected:Npn \__web_picture_number:nn #1#2 {
    \bool_set_false:N \l__web_picture_back_bool
    \tl_clear:N \l__web_picture_tag_tl
    \__web_picture_tag:w #2 \tag \q_stop
    % (a \tag*{} gives an empty tag, and no number)
    \bool_if:NF \l__web_picture_back_bool {
        \clist_if_in:NnT \c__web_numbered_clist {#1} {
            \tl_if_in:nnF {#2} {\notag} {
                \tl_if_in:nnF {#2} {\nonumber} {
                    \refstepcounter{equation}
                    \tl_set:Nn \l__web_picture_tag_tl { (\theequation) }
                }
            }
        }
    }
}
% the picture, its number, and the box of a \qedhere, in the group of the
% number's step
\cs_new_protected:Npn \__web_picture_set:n #1 {
    \bool_gset_false:N \g__web_picture_qed_bool
    \cs_set_eq:NN \__web_picture_qedhere: \qedhere
    \cs_set_protected:Npn \qedhere { \bool_gset_true:N \g__web_picture_qed_bool }
    \cs_set_protected:Npn \tag { \@ifstar \use_none:n \use_none:n }
    \cs_set_eq:NN \notag \prg_do_nothing:
    \cs_set_eq:NN \nonumber \prg_do_nothing:
    \__web_picture_body:n {#1}
    \exp_args:Ne \__web_picture:n { \exp_not:N \hbox { \exp_not:V \l__web_picture_body_tl } }
    \__web_picture_punct:
    \bool_if:NT \l__web_picture_back_bool
        { \setcounter{equation}{\l__web_picture_count_int} }
    % (in the upright font, as amsmath and LaTeX set a number)
    \tl_if_empty:NF \l__web_picture_tag_tl {
        \HCode{<span~class="tex8web-picture-tag">}
        \hbox { \normalfont \normalcolor \l__web_picture_tag_tl }
        \HCode{</span>}
    }
    \bool_if:NT \g__web_picture_qed_bool { \hbox { \__web_picture_qedhere: } }
}
% \tag{text} or \tag*{text} in the body: the tag, which a \label names
\cs_new_protected:Npn \__web_picture_tag:w #1 \tag #2 \q_stop {
    \tl_if_blank:nF {#2} { \__web_picture_tag:nw #2 \q_stop }
}
\cs_new_protected:Npn \__web_picture_tag:nw #1#2 \q_stop {
    \tl_if_eq:nnTF {#1} {*}
        { \__web_picture_tag_star:nw #2 \q_stop }
        { \tl_set:Nn \l__web_picture_tag_tl { (#1) } \__web_picture_label:n {#1} }
}
\cs_new_protected:Npn \__web_picture_tag_star:nw #1#2 \q_stop {
    \tl_set:Nn \l__web_picture_tag_tl {#1}
    \__web_picture_label:n {#1}
}
% (a step of the counter, with the tag as its number, for the anchor
% before the picture that a \label names, wherever it is; the counter is
% set back after the picture's \label. hyperref names the anchor from
% \theHequation, which is given a name of its own, so that the next
% numbered equation's anchor has another)
\int_new:N \l__web_picture_count_int
\int_new:N \g__web_picture_tag_int
\bool_new:N \l__web_picture_back_bool
\cs_new_protected:Npn \__web_picture_label:n #1 {
    \int_set:Nn \l__web_picture_count_int { \value{equation} }
    \bool_set_true:N \l__web_picture_back_bool
    \cs_set:Npn \theequation {#1}
    \int_gincr:N \g__web_picture_tag_int
    \cs_set:Npe \theHequation { tag. \int_use:N \g__web_picture_tag_int }
    \refstepcounter{equation}
}
% (LaTeX's eqnarray takes no \tag)
\cs_new_protected:Npn \__web_if_env_picture:nnTF #1#2#3#4 {
    \str_if_in:nnTF {#1} {eqnarray}
        { \__web_if_bare_picture:nTF {#2} {#3} {#4} } { \__web_if_picture:nTF {#2} {#3} {#4} }
}
% (amsmath's flalign, for the glyphs of a saved box: tex4ht's mathjax
% option names it flalign[equation], so \RestoreMathJaxEnvironments
% leaves tex8web's in place, mathjax-latex-4ht.4ht)
\cs_set_eq:NN \__web_saved_flalign: \flalign
\cs_set_eq:NN \__web_saved_endflalign: \endflalign
\clist_map_inline:nn { equation, equation*, align, align*, gather, gather*, alignat, alignat*,
        flalign, flalign*, multline, multline*, eqnarray, eqnarray* } {
    \cs_if_exist:cT {#1} {
        \RenewDocumentEnvironment{#1}{+!b}{
            \__web_if_env_picture:nnTF {#1} {##1}
                { \__web_picture_display:nn {#1} {##1} }
                { \__web_saved_env:nn {#1} {##1} \__web_display:nn {#1} {##1} }
        }{}
    }
}
\RenewDocumentEnvironment{displaymath}{+!b}{
    \__web_if_bare_picture:nTF {#1}
        { \__web_picture_display:nn {displaymath} {#1} }
        { \__web_saved_env:nn {displaymath} {#1} \__web_mathjax:nn {#1} {displaymath} }
}{}

% Math in a saved box (\sbox, \savebox, lrbox). tex4ht's mathjax option
% writes a formula's source as specials and typesets nothing, and tex4ht
% drops these specials in a picture. So a box placed in a TikZ node would
% lose its math. In such a box each formula is also typeset as TeX would,
% as in a picture. dvisvgm draws these glyphs in a picture. Elsewhere
% tex4ht ignores them, as they lie between the specials t4ht@[ and t4ht@]
% (info4ht.4ht). Their own specials are dropped, so tex4ht's math hooks
% write no HTML. So are their labels, as the source writes them: tex4ht's
% \label writes to the aux file only if \if@filesw (latex.4ht). The
% equation counter is set back, so the source numbers its rows as if the
% glyphs were not there, and so is amsthm's QED stack, which a \qedhere
% changes globally. (\ht:special and \m:op, tex4ht.sty; \nested:math,
% html4-math.4ht)
\bool_new:N \l__web_saved_bool
\int_new:N \l__web_saved_int
\tl_new:N \l__web_saved_tl
\tl_new:N \l__web_saved_qed_tl
\regex_const:Nn \c__web_saved_amp_regex { \cL\& }
\web@needs{n:sbox:}
\web@needs{o:sbox:}
\web@needs{gobble:fbox}
\web@needs{gobble:frame}
\web@needs{oncenter:list}
\web@needs{this:listConfigure}
\web@needs{SaveEndP}
\web@needs{a:MathjaxBlock}
\web@needs{b:MathjaxBlock}
\web@needs{:HandleMathjaxCatcodes}
\web@needs{c:display}
\web@needs{ht:special}
\web@needs{nested:math}
\web@needs{RestoreMathJaxEnvironments}
% tex4ht's \everydisplay ends a $$ display at once and starts another
% (\NewConfigure{$$}, tex4ht.sty), so a saved box would hold an empty
% display above the glyphs' display, which the PDF does not have. In a
% saved box, the glyphs and the source are set in the first display.
%
% A TikZ node draws a saved box as TeX set it, but the box is set outside
% a picture. There latex.4ht's \centering sets \leftskip to zero, and its
% \fbox and \framebox drop their frames. In a saved box they are set as in
% LaTeX, and so are the frame's rules, by \__web_saved_frame:n below.
\cs_new_protected:Npn \__web_saved_start: {
    \bool_set_true:N \l__web_saved_bool
    \tex_everydisplay:D { \if_meaning:w \EndPicture \:UnDef \__web_saved_display: \fi: }
    \tl_put_right:Nn \centering { \leftskip \@flushglue \parfillskip \c_zero_skip }
    \cs_set_eq:cN {gobble:fbox} \__web_saved_fbox:nn
    \cs_set_eq:cN {gobble:frame} \__web_saved_framebox:w
    \cs_set_eq:NN \@frameb@x \__web_saved_frame:n
    \regex_replace_once:nnN { \c{item} (.*) \c{par} \Z }
        { \c{__web_saved_depth:} \c{item} \1 } \center
    % (once, in a box saved in another)
    \cs_if_eq:cNF {oncenter:list} \__web_saved_center: {
        \cs_set_eq:Nc \__web_saved_center_list: {oncenter:list}
        \cs_set_eq:cN {oncenter:list} \__web_saved_center:
    }
}
\cs_new_protected:Npn \__web_saved_center: {
    \__web_saved_center_list:
    \tl_put_right:Nn \this:listConfigure { \cs_gset_eq:NN \:DT \__web_saved_item:w }
}
\cs_new:Npn \__web_saved_display_end: { }
\cs_new_protected:Npn \__web_saved_display: {
    \use:c {a:MathjaxBlock} \use:c {:HandleMathjaxCatcodes}
    \tex_everymath:D {} \tex_everydisplay:D {}
    % (the tag that a:MathjaxBlock chose, which the display's group ends)
    \cs_gset:Npe \__web_saved_display_end: {
        \cs_set:Npn \exp_not:N \web@mathtag { \web@mathtag } \exp_not:c {b:MathjaxBlock}
    }
    \group_insert_after:N \__web_saved_display_end:
    \use:c {c:display}
}
% LaTeX's \fbox and \framebox, where tex4ht writes no space for the
% frame's sides and \fboxsep: only the text is set outside t4ht@( and
% t4ht@), and t4ht@[ and t4ht@] (info4ht.4ht). The \hbox that latex.4ht
% sets around the text of an \fbox is dropped, as tex4ht writes a space
% at the move past its end. #1: LaTeX's command.
\cs_new_protected:Npn \__web_saved_fbox:nn #1#2 {
    \use:c {ht:special} {t4ht@(}
    #1 { \use:c {ht:special} {t4ht@)} \use_ii:nn #2 \use:c {ht:special} {t4ht@[} }
    \use:c {ht:special} {t4ht@]}
}
\cs_new_protected:Npn \__web_saved_framebox:w #1 [#2] [#3] #4 {
    \use:c {ht:special} {t4ht@(}
    #1 [#2] [#3] { \use:c {ht:special} {t4ht@)} #4 \use:c {ht:special} {t4ht@[} }
    \use:c {ht:special} {t4ht@]}
}
% LaTeX's \@frameb@x, with the frame's rules first and between t4ht@[ and
% t4ht@], where tex4ht writes no rule (info4ht.4ht). The text stays on
% the line's baseline, as tex4ht starts a new line where a glyph sits
% lower than the one before. #1: a kern inside each side.
\cs_new_protected:Npn \__web_saved_frame:n #1 {
    \hbox_set:Nn \@tempboxa { \tex_kern:D \fboxrule #1 \box_use_drop:N \@tempboxa #1 \tex_kern:D \fboxrule }
    \dim_set:Nn \@tempdima { \box_ht:N \@tempboxa + \fboxsep + \fboxrule }
    \dim_set:Nn \@tempdimb { \box_dp:N \@tempboxa + \fboxsep + \fboxrule }
    \box_set_ht:Nn \@tempboxa { \@tempdima }
    \box_set_dp:Nn \@tempboxa { \@tempdimb }
    \hbox:n {
        \use:c {ht:special} {t4ht@[}
        \tex_vrule:D width \box_wd:N \@tempboxa height \@tempdima
            depth \dim_eval:n { \fboxrule - \@tempdima }
        \tex_kern:D - \box_wd:N \@tempboxa
        \tex_vrule:D width \box_wd:N \@tempboxa height \dim_eval:n { \fboxrule - \@tempdimb }
            depth \@tempdimb
        \tex_kern:D - \box_wd:N \@tempboxa
        \tex_vrule:D width \fboxrule height \@tempdima depth \@tempdimb
        \tex_kern:D \dim_eval:n { \box_wd:N \@tempboxa - 2 \fboxrule }
        \tex_vrule:D width \fboxrule height \@tempdima depth \@tempdimb
        \tex_kern:D - \box_wd:N \@tempboxa
        \use:c {ht:special} {t4ht@]}
        \box_use_drop:N \@tempboxa
    }
}
% A center's first line in a saved box. html4.4ht's \center ends the
% paragraph that \item starts, and its list hook ends the next one. Then
% latex.4ht's \item, at \:DT, writes \hfill\break. These leave two empty
% lines above the text, which the PDF does not have. In a saved box the
% first, which holds the center's <div>, is kept with no height, and the
% text starts the second, as far below the line above as in the PDF.
% (\:DLL and \:DT, latex.4ht)
\dim_new:N \l__web_saved_depth_dim
\cs_new_protected:Npn \__web_saved_depth: { \dim_set:Nn \l__web_saved_depth_dim { \tex_prevdepth:D } }
\cs_new_protected:Npn \__web_saved_item:w #1#2 {
    \tex_setbox:D 0 \tex_lastbox:D
    \dim_compare:nNnF \l__web_saved_depth_dim < { -999pt } { \tex_unskip:D }
    \tex_prevdepth:D -1000pt
    \tex_box:D 0
    \tex_prevdepth:D \l__web_saved_depth_dim
    \mode_leave_vertical:
}
% The end of the paragraph before a box saved with \sbox or \savebox is
% kept out of the box, as html4.4ht keeps it out of an lrbox.
\cs_set_protected:cpn {n:sbox:} #1#2 { \use:c {o:sbox:} {#1} { { \SaveEndP \__web_saved_start: #2 \RecallEndP } } }
\Configure{lrbox}{\__web_saved_start:}
% #1: the formula as LaTeX source. tex4ht reads & in \(, \[, and $ as a
% letter (\:HandleMathjaxCatcodes, html4-math.4ht)
\cs_new_protected:Npn \__web_saved:n #1 {
    \bool_if:NT \l__web_saved_bool {
        \int_set:Nn \l__web_saved_int { \value{equation} }
        \cs_if_exist:NT \QED@stack { \tl_set_eq:NN \l__web_saved_qed_tl \QED@stack }
        \use:c {ht:special} {t4ht@[}
        \group_begin:
            \cs_set_eq:cN {ht:special} \use_none:n
            \@fileswfalse
            \bool_set_false:N \l__web_saved_bool
            \cs_set:cpn {EndPicture} {}
            \use:c {nested:math}
            \RestoreMathJaxEnvironments
            \cs_set_eq:NN \flalign \__web_saved_flalign:
            \cs_set_eq:NN \endflalign \__web_saved_endflalign:
            % (siunitx's math, which tex4ht and tex8web pass to MathJax,
            % set by siunitx: see \__web_saved_si_math:n)
            \cs_if_exist:NT \__web_saved_si_math:n
                { \cs_set_eq:NN \__siunitx_print_math_auxii:n \__web_saved_si_math:n }
            \tl_set:Nn \l__web_saved_tl {#1}
            \regex_replace_all:NnN \c__web_saved_amp_regex { \cT\& } \l__web_saved_tl
            \l__web_saved_tl
        \group_end:
        \use:c {ht:special} {t4ht@]}
        \int_gset_eq:NN \c@equation \l__web_saved_int
        \cs_if_exist:NT \QED@stack { \tl_gset_eq:NN \QED@stack \l__web_saved_qed_tl }
    }
}
% An environment's glyphs: its \end sets \@ignoretrue, as amsmath's and
% LaTeX's displays do, and the source's \end then skips the space after
% it, as in the PDF. Otherwise the glyphs' own \end would use the flag
% up inside the group (\end, ltmiscen.dtx). tex4ht's \end writes \global
% before it (latex.4ht), so it stays an assignment, one that changes
% nothing.
\cs_new_protected:Npn \__web_saved_env:nn #1#2 {
    \__web_saved:n { \cs_set:Npn \@ignorefalse { \tex_let:D \if@ignore \if@ignore } \begin{#1} #2 \end{#1} }
}
% A \[ \] with a \qedhere: \nested:math's \[ is a bare $$, in which
% amsthm warns that \qedhere may not work and sets the box a \quad after
% the formula (\qedhere, amsthm.sty). amsmath's \[ is equation*, which
% puts the box at the right margin, as in the PDF. Its \end clears the
% flag it sets, as the source's \] is no environment
\cs_new_protected:Npn \__web_saved_qed:n #1 {
    \cs_if_exist:cTF {equation*}
        { \__web_saved:n { \begin{equation*} #1 \end{equation*} } }
        { \__web_saved:n { \[#1\] } }
}
\web@needs{AltMathOne}
\web@needs{AltlMath}
\cs_set_eq:NN \__web_saved_AltMathOne:w \AltMathOne
\cs_set_eq:NN \__web_saved_AltlMath:w \AltlMath
% (in math already, after $. \[ \] and $$ call \__web_saved:n above,
% where they are not pictures)
\cs_set:Npn \AltMathOne #1 $ { \__web_saved:n {#1} \__web_saved_AltMathOne:w #1 $ }
\cs_set:Npn \AltlMath #1 \) { \__web_saved:n { \(#1\) } \__web_saved_AltlMath:w #1 \) }
\ExplSyntaxOff

\makeatother
